\section{Audit, integration, and further research}
\label{sec:research}

\subsection{Dependency and claim audit}

The proof is finite and self-contained apart from standard real and
integer arithmetic and the finiteness of the usual van der Waerden
numbers used to name \(W_r(k)\). The lower-bound construction itself
does not use a quantitative upper-bound theorem.
The indispensable chain is:
\begin{enumerate}
 \item \Cref{lem:prime-estimates,lem:dilation,lem:mesh} give the
       cyclic coordinates, a controlled dilation, and stable mesh
       neighborhoods.
 \item \Cref{lem:global,lem:local} bound global words and local
       patterns through a specified literal label.
 \item \Cref{prop:outer-tests} applies a finite local lemma to
       distinct-label rows, residual light sets, and local patterns.
 \item \Cref{lem:return,lem:short,thm:dichotomy} convert repeated
       labels into either a global affine lift or near-uniformity.
 \item \Cref{lem:multiplicity,lem:signature,thm:abstract} control
       repeated keys and both kinds of bad event.
 \item \Cref{thm:cyclic} checks every parameter loss, and
       \cref{prop:product,prop:periodic,prop:palette} supply the
       uniform all-color conclusions.
\end{enumerate}

Several mathematical boundaries deserve explicit recording.
The key-count argument uses separation of the actual \(q\)-adic
embedding, not a generic assertion about thin annuli.
The affine event count preserves key identities; a count only up to
renaming keys would be insufficient.
The local lemma uses independent variables attached to labels, while
the final perturbation uses a separate independent family attached to
keys. Independence is never asserted between progressions.
The three-color result controls distance from a constant word and
does not imply occurrence of all three colors.
Finally, the mixed-radix product works even when its moduli share
prime factors, because it uses a divisibility chain, not the Chinese
remainder theorem.

\begin{center}
\small
\begin{tabular}{@{}>{\raggedright\arraybackslash}p{.29\textwidth}>{\raggedright\arraybackslash}p{.27\textwidth}>{\raggedright\arraybackslash}p{.34\textwidth}@{}}
\toprule
Quantity&Displayed source choice&Choice or bound here\\
\midrule
Dimension&\(\cl{k^{1/10}}\)&\(\cl{k^{1/3}(\log k)^{-2/3}}\)\\
Small-prime cutoff&\(D^2\)&\(\cl{AD(\log k)^3}\)\\
Outer binary balance&at least \(k/100\) of each&at least
\((1/2-4/\log k)k\) of each, or affine\\
Key multiplicity&at most four&at most two\\
Perturbation probability&\(k^{-1/20}\)&
\(Bk^{-1/3}(\log k)^{5/3}\)\\
Binary \(k\log k\) coefficient&\(10^{-5}\)&\(1/24\)\\
Color products&binary digits&optimized binary and ternary digits\\
Distance guarantee&nonmonochromatic&order
\(k^{2/3}(\log k)^{5/3}\)\\
\bottomrule
\end{tabular}
\end{center}

\subsection{What the companion computation establishes}

The accompanying \texttt{verify.py} uses exact rational arithmetic for
the palette and probability checks. It verifies the two coefficient
cancellations in \eqref{eq:cancellations}, the correction polynomials
in \eqref{eq:reversion}, the palette optimization over a finite
validation range, and exact mixed-radix floor identities for small
noncoprime and mixed bases. It also checks finite increment spaces
against the rich-event estimate with repeated keys, including
incompatible prescriptions. Its report records the ranges and counts.

These checks are useful for catching algebra, ceiling, carry, and
shared-variable errors. They do not numerically construct the enormous
colorings at the eventual threshold or validate an infinite theorem
by enumeration. The proofs above, including the asymptotic domination
estimates and the local lemma, are the mathematical justification.
The figure is generated from the exact linear inequalities
\eqref{eq:power-ledger}, not from experimental estimates of
van der Waerden numbers.

\subsection{Suggested placement and formalization route}

This package is intended for a distinct research subdirectory under
\texttt{Combinatorics/}\allowbreak\texttt{Ramsey/}\allowbreak\texttt{Research/},
since its new claims are lower
bounds and robustness properties, rather than exact replacements of
the Gowers upper-bound interfaces.
It need not alter the existing 120-statement Gowers ledger.
The source citations are pinned in \texttt{SOURCE\_PROVENANCE.json};
the new article, bibliography, verification program, figure, and build
instructions form a standalone package.

A formalization can be staged at precise mathematical boundaries:
\begin{enumerate}
 \item Prove the norm-band lemma for finite separated subsets of a
       real affine line and then apply the cyclic coordinate
       separation lemma.
 \item Formalize \cref{thm:abstract} for a finite indexed family of
       configurations. This part is independent of the torus meshes.
 \item Formalize the mixed-radix transfer and the exact integer
       palette formula. These are likewise independent of the
       asymptotic analytic estimates.
 \item Formalize mesh-width and neighborhood estimates, arrangement
       sign counts including equality, and the amended residual
       event family.
 \item Assemble the return dichotomy and the endpoint domination
       inequalities, with constants fixed in the stated order.
\end{enumerate}
The new statements should acquire formal status only after those proofs
are actually checked in the selected prover. The existing source's
formalized coefficient \(10^{-5}\) does not transfer formal status to
coefficient \(1/24\) by citation.

\subsection{Further research questions}

\begin{question}[Reduce global label entropy]\label{q:global}
Can the actual global words of the \(q\)-adic embedding be counted
more efficiently than by treating all scalar starts and steps
independently?
\end{question}

The proof currently pays \(D\log\lambda\), of size
\(D^2\) times logarithmic factors.
The scalar homomorphisms are strongly dependent, so their independent
product count may be wasteful.
For orientation, if a valid replacement gave a power-scale cost
\(D^u\), while leaving the quadratic cost \(k/D\) and other steps
intact, the same exponent calculation would permit a binary
coefficient approaching \(1/[8(u+1)]\).
This conditional calculation identifies what an improved count would
buy; it does not supply the count.
Conditioning the word count on a closest return or on a common
carry pattern is a concrete first direction.

\begin{question}[Exploit repeated-key structure]\label{q:keys}
Do nonaffine progressions with many double keys necessarily have
smaller signature complexity, or many incompatible outer-color
prescriptions?
\end{question}

The general two-point geometric bound is best possible, but the final
union bound uses it in every rich progression.
A double key requires the corresponding line segment to meet opposite
sides of a norm minimum. A structural theorem making extensive pairing
either compressible or incompatible with the desired target color
could increase the forced-increment count without changing the mesh.
Proving an average statement is insufficient unless it controls all
remaining bad events.

\begin{question}[Improve the logarithmic correction]\label{q:log}
Can the coefficient \(5/24\) of \(k\log\log k\) in the binary bound
be reduced while retaining the leading coefficient \(1/24\)?
\end{question}

The present value comes from the anchored count, a shrinking balance
tolerance, and the two signature bounds.
A sharper treatment of local dependencies or an entropy bound using
the distribution of anchor widths may change it.
The optimization discussion establishes neither a lower bound on
actual entropy nor a universal second-order obstruction.

\begin{question}[A larger robustness scale]\label{q:robust}
Can one obtain \(N\ge\exp(c k\log k)\), with fixed \(c>0\), while
forcing every \(k\)-term progression to contain a fixed positive
fraction of both colors?
\end{question}

The fixed-power result \cref{thm:tradeoff} reaches every
\(k^\theta\), \(2/3<\theta<1\), at the price of a smaller coefficient,
but its permitted \(c\) tends to zero as \(\theta\to1\).
The endpoint result attains the largest coefficient supported by the
current ledger and a logarithmic enhancement of \(k^{2/3}\).
Whether this tradeoff reflects the problem or just the perturbation
method is not settled by the report.

\begin{question}[Construct and certify the succinct tables]\label{q:algorithm}
Can successful tables from \cref{cor:description} be found and
certified in time polynomial in their description length?
\end{question}

An outer local-lemma algorithm is only part of the task.
The eligible patterns must be enumerated or separated effectively,
and the rich progression failures in the second stage must also be
controlled.
An acceptable certificate should establish coverage of every
configuration rather than merely test a large random sample.
The algebraic form of the mesh endpoints makes exact comparisons
possible, so the problem is combinatorial and quantitative, not an
obstruction from noncomputable real parameters.

\begin{question}[Explicit onset and finite constructions]\label{q:onset}
What explicit numerical constants and onset can be extracted, and
which simplifications help at feasible values of \(k\)?
\end{question}

The current proof deliberately separates existence from numerical
optimization. Tracking an explicit arrangement constant, local
incident weight, and prime-count constant would give an effective
but potentially enormous threshold.
Such a calculation should be compared with direct finite coloring
constructions; the large asymptotic coefficient alone is not evidence
of a better small-\(k\) record.

\begin{question}[Asymmetric progression lengths]\label{q:offdiag}
Can the balanced-or-affine construction be optimized when different
colors forbid different progression lengths?
\end{question}

Biased outer color probabilities and nonuniform increment probabilities
would lead to a multiobjective version of
\cref{thm:abstract}.
The central issue is that the geometric period and mesh parameters
would need to handle several lengths simultaneously.
An exact asymmetric optimization should distinguish a color's
extinction event from proximity to a constant word, a distinction
already essential in the ternary case here.

\begin{question}[Other arithmetic configuration families]\label{q:patterns}
Which systems beyond linear progressions admit simultaneously
bounded key multiplicity and a small literal exceptional-signature
space?
\end{question}

\Cref{thm:abstract} already applies to arbitrary finite configuration
families once those two inputs are supplied.
For polynomial configurations a norm evaluated along the lifted path
has higher degree, but a degree bound by itself does not supply the
separation geometry used in \cref{lem:multiplicity}.
A valid extension must prove the new geometric multiplicity lemma,
rather than assume that the line argument transfers.

\begin{question}[Compare small- and large-palette methods]\label{q:palettes}
Can the present small-palette construction be combined with the
Fox--Hunter shifted-product method to improve the transition where
the number of colors is polylogarithmic in \(k\)?
\end{question}

The direct binary--ternary product loses efficiency as a function of
\(\log r\), whereas Fox--Hunter is substantially stronger once its
large-palette hypothesis holds.
A useful combination would need to preserve quantitative defect
information through each shifted product, with every density loss
recorded. The robust cyclic theorem supplies such an input, but no
improvement in that transition regime is claimed here.

\begin{question}[Kernel-checked constants and proof interfaces]
\label{q:formal}
Can the strengthened finite interfaces, their asymptotic assembly,
and the explicit exponent \(1/24\) be formalized without importing
the stronger endpoint as an assumption?
\end{question}

The return argument, residual-set definition, and shared-key random
variables make useful audit targets.
Their exact quantifiers are more informative than a theorem wrapper
asserting only the final existence of constants.
A successful formalization would both verify this report and provide
reusable lower-bound interfaces beside ProveIt's existing
upper-bound formalization program.

\subsection{Scope of the outcome}

The report proves a substantial quantitative refinement of the selected
repository construction, with an explicit leading coefficient,
logarithmic correction, robustness guarantee, inverse form, and exact
palette optimization.
It does not determine the true asymptotic growth of \(W_2(k)\), prove
optimality of the new constants, claim the best estimate for every
color regime, or provide a formalized proof of the strengthened
theorems. Those boundaries specify what the next research and
verification steps must address.
